\documentclass[10pt]{article}
\usepackage{graphicx} 
\usepackage{amsmath}
\usepackage{amsthm}
\usepackage{titlesec}

\usepackage{fontawesome5}

\usepackage[top=25mm,bottom=25mm,left=20mm,right=20mm
]{geometry}

\titleformat{\section}{\large\bfseries}{}{0em}{}[\titlerule]

\usepackage[dvipsnames]{xcolor}
\usepackage{hyperref}
\hypersetup{
    colorlinks = true,
    urlcolor  = {blue!50!black}, 
}


\begin{document}

\pagestyle{empty}
\begin{center}
    {\textbf{\huge Milène Mandeville}}
    
    \bigbreak
    \Large{Applied Mathematics Student} \normalsize{(She/Her)}\vspace{3mm}
\end{center}
\begin{center}
    \faIcon{phone} \href{https://wa.me/33668298113}{+33 6 68 29 81 13} \hspace{1cm}
    \faIcon{envelope} \href{mailto:milenemandeville@yahoo.com}{milenemandeville@yahoo.com} \hspace{1cm}
    \faIcon{linkedin} \href{https://www.linkedin.com/in/milènemandeville/}{milènemandeville}\vspace{3mm}
\end{center}
\begin{minipage}[t]{0.45\linewidth}
\section*{Languages}
\begin{minipage}[t]{0.5\linewidth}
\begin{itemize}
    \item French (native)
    \item English (C1)
\end{itemize}
\end{minipage}\hfill 
\begin{minipage}[t]{0.5\linewidth}
    \begin{itemize}
    \item Dutch (B1)
    \item German (B1)
\end{itemize}
\end{minipage}
\end{minipage}\hfill
\begin{minipage}[t]{0.5\linewidth}
\section*{Programming skills \quad\, \normalfont{\normalsize{\faGithub\, \href{https://github.com/MileneMdvl}{MileneMdvl}}}}
\begin{minipage}[t]{0.5\linewidth}
\begin{itemize}
    \item Python (advanced)
    \item LaTex (advanced)
    \item Julia (advanced)
\end{itemize}
\end{minipage}\hfill
\begin{minipage}[t]{0.5\linewidth}
\begin{itemize}
    \item Matlab (intermediate)
    \item Fortran (intermediate)
\end{itemize}
\end{minipage}
\end{minipage}\hfill
% \begin{minipage}[t]{0.3\linewidth}
% \section*{Diplomas}
% \begin{itemize}
%     \item BSc in Applied \\ Mathematics (2024)
%     \item French Baccalaureate in Science -- honours (2020)
% \end{itemize}
% \end{minipage}\hfill 


\section*{Education}
\begin{itemize}
    \item \textbf{MSc Applied Mathematics}, Rijksuniversiteit Groningen
    (2024--2026) 

    \textbf{Relevant coursework}: Computational Fluid Dynamics, Finite Element
    Methods, Coupled Problems, Inverse Problems in Imaging, Numerical Simulation
    of Turbulence, Randomised Numerical Linear Algebra
    \item \textbf{BSc Applied Mathematics}, Rijksuniversiteit Groningen
    (2020--2024) 
 
    % \item Agora High School, Puteaux France (2017-2020)
\end{itemize}
\section*{Research}
\begin{itemize}
    \item \textbf{Master Thesis}: University of Groningen (2026) \begin{itemize}
        \item \textbf{Title}: \href{https://milenemdvl.github.io/CV/theses/MasterThesis.pdf}{Unstructured Finite Volume Discretisation Scheme for Turbulent Navier-Stokes}
        \item \textbf{Supervisors}: prof. dr. ir. R.W.C.P Verstappen and dr. ir. R. Luppes 
        \item \textbf{Description}: In this paper, I present a Finite Volume discretisation scheme for two- and three-dimensional unstructured grids obtained with Delaunay Triangulation. This triangulation is chosen such that the vertices only lie on the boundaries of the triangles' circumcircles. I first go over the basics of turbulence modelling, such as Large Eddy Simulation and filtering. With the aim of the divergence and gradient being adjoints of one another, I present the derivations for the discrete operators, and I give an example for a simple two-dimensional grid. I solve the pressure via a Poisson problem where the intermediate linear systems are solved iteratively. I verify the code using the method of manufactured solutions, in which I present velocity and density solutions as linear combinations of sines and cosines and aim to recover these solutions.
    \end{itemize}
    \item \textbf{Internship}: Von Karman Insitute for Fluid Dynamics (2026) \begin{itemize}
        \item \textbf{Title}: \href{https://milenemdvl.github.io/CV/theses/InternshipThesis.pdf}{Towards Efficient Galerkin Projection in
        Lagrangian ROMs for Sloshing}
        \item \textbf{Supervisors}: prof. M. A. Mendez, ir. S. Ahizi,  dr. J. Koellermeier, and dr. ir. R.
        Luppes,
        \item \textbf{Description}: In this paper I derive a reduced order model for
        smoothed particle hydrodynamics, by rewriting the standard neighbour
        sums as dense matrix operations, and using proper orthorgonal
        decomposition and Galerkin projection to obtain a linear time invariant
        model. However, the kernel transformations inherent to particle dynamics
        introduce a nonlinear state dependent term. To this end, I evaluate
        three methods: linearisation, random Fourier features, and a singular
        value decomposition based interpolation. 
    \end{itemize}
    \item \textbf{Bachelor Thesis}: University of Groningen (2023) \begin{itemize}
        \item \textbf{Title}: \href{https://milenemdvl.github.io/CV/theses/BachelorThesis.pdf}{Average Surface Roughness for 2-Dimensional Atmospheric Flow}
        \item \textbf{Supervisors}: dr. ir. R. Luppes and dr. A. E. Sterk 
        \item \textbf{Description}: In this paper I present a study of four different methods of computing an
        average (effective) roughness length (following André and Blondin (1986), Taylor
        (1989), de Vries et. al. (2003), and taking the average of the
        logarithmic law).
    \end{itemize}
\end{itemize}

\section*{Professional experience}
\begin{itemize}
\item \textbf{Acadomia (tutoring company): online teacher} (2024--2026)
\begin{itemize}
    \item Teaching mathematics to middle school and high school students in French
\end{itemize}
\item \textbf{University of Groningen: teaching assistant} (2023--2026)\\ 
    For the following Bachelor courses, I gave tutorials and/or computer labs in
    Python. I
    also graded homeworks and exams, and evaluated students
    with oral examinations for the courses with labs. \begin{itemize}
    \item Mathematical Tools for BME (2025) -- C. Karakurt, PhD
    \item Project Mathematical Modelling (2025) -- prof. dr. ir. R. Verstappen
    \item Computational Methods of Science (2024--2026) -- prof. dr. ir. R. Verstappen, dr. ir. R. Luppes
    \item Mathematical Modelling (2023--2026) -- dr. S. Trenn, prof dr. ir. R. Verstappen
    \item Advanced Systems Theory (2023) -- prof. M.K. Camlibel
\end{itemize}
    % \item \textbf{Puteaux townhall Public Library (Paris La Défense suburb): summer job} (2023) \begin{itemize}
%     \item Welcomed and helped visitors
%     \item Sorted books using Orphée software
% \end{itemize}
% \item \textbf{Penoyre \& Prasad Architectes: intern} (2017)
% \begin{itemize}
%     \item Observational internship
% \end{itemize}
\end{itemize}


\section*{Other experience}
\begin{itemize}
    \item \textbf{LaTex Workshop} (2025--2026)\begin{itemize}
        \item Teaching first year Bachelor Students how to use LaTex 
        \item Focusing on equations, theorems, figures, tables, BibTex
        management and Beamer 
        \item Introduction to TikZ
    \end{itemize}
    \item \textbf{Secretary of a committee of the FMF (study association)} (2023--2026)\begin{itemize}
    \item Writing minutes summarising the meetings 
    \item Writing emails inviting professors to give a talk 
    \item Organising the annual integration competition and a series of informal lectures
\end{itemize}
\end{itemize}

\section*{References}
\begin{itemize}
    \item Prof. dr. ir. Roel Verstappen, University of Groningen, \href{mailto:r.w.c.p.verstappen@rug.nl}{r.w.c.p.verstappen@rug.nl}
    \item Dr. ir. Roel Luppes, University of Groningen, \href{mailto:r.luppes@rug.nl}{r.luppes@rug.nl}
    \item Prof. Miguel Mendez, von Karman Institute for Fluid Dynamics, 
    \href{mailto:miguel.alfonso.mendez@vki.ac.be}{miguel.alfonso.mendez@vki.ac.be}
\end{itemize}

\qed

\end{document}
